自定义的损失写完，训练跑起来，曲线在降。降得对不对，这时候没有任何东西会告诉你。梯度里少乘一个 \(1/n\)、转置放在了错误的一侧、bias 的梯度忘了沿 batch 轴求和——这几类错误都不抛异常，它们只是让曲线比应有的样子平一点、慢一点。框架也帮不上忙，因为它保证的是``按你写的前向求导''，不是``你写的前向是你想要的那个函数''。

本单元把前三个单元的工具用到真实模型上：形状先声明，微分整理成 \(dL=\trace(G^{\trans}dX)\)，梯度从这个形状里直接读出来。三篇按复杂度递增，方法只有一套。

\begin{center}
\begin{tikzpicture}[x=1cm,y=1cm,font=\small,>=stealth]
  \node[font=\footnotesize,black!65] at (4.4,3.2) {同一套读取程序：把 \(dL\) 整理成 \(\trace(G^{\trans}dX)\)，读出 \(\nabla_XL=G\)};
  \draw[black!25] (-1.8,2.9) -- (10.6,2.9);
  \foreach \x/\ta/\tb in {0/{回归与分类}/{\(X^{\trans}(p-y)\)},4.4/{softmax 配交叉熵}/{\(p-y\)},8.8/{计算图}/{局部导数 \(\times\) 顺序}}
    {\draw[draw=black!55,fill=blue!7,rounded corners=2pt] (\x-1.7,1.35) rectangle (\x+1.7,2.45);
     \node[font=\footnotesize] at (\x,2.12) {\ta};
     \node[font=\footnotesize] at (\x,1.68) {\tb};}
  \draw[->,black!55] (1.75,1.9) -- (2.65,1.9);
  \draw[->,black!55] (6.15,1.9) -- (7.05,1.9);
  \node[font=\footnotesize,black!60] at (0,0.85) {出错先查形状与 \(1/n\)};
  \node[font=\footnotesize,black!60] at (4.4,0.85) {出错先查中间量与溢出};
  \node[font=\footnotesize,black!60] at (8.8,0.85) {出错先查顺序与广播轴};
  \draw[->,red!60!black,line width=1.1pt] (-1.7,0.1) -- (10.5,0.1);
  \node[font=\footnotesize,red!55!black,anchor=south] at (4.4,0.18) {手推一次，是为了知道框架在算什么};
\end{tikzpicture}

\emph{三篇是同一套推导落在三处：一层模型上的手推、一次配对化简、整张图上的自动执行；下面一行是它们共同的用途。}
\end{center}

第一篇在同一副线性预测骨架上比较两种损失。平方损失和对数似然的梯度都收缩成 \(X^{\trans}\delta\) 的样子，差别只在 \(\delta\) 那一格，\(X^{\trans}\) 那一格原样不动。这个共同结构不是巧合，它是后面所有推导的模板。

第二篇处理一个看上去必须硬算的对象。softmax 的 Jacobian 是 \(K\times K\) 的，按定义走要先把它物化出来，\(K=1000\) 时每个样本一百万个数；而把损失写成 logits 的函数之后，指数在取对数那一步就消掉了，梯度只剩 \(p-y\)。化简的来源是 softmax 与交叉熵配成一对，这也顺带解释了框架为什么要提供一个吃 logits 的合并算子，以及为什么它同时更省更稳。

第三篇把手推换成机械执行。计算图上每个节点只需要知道自己的局部导数，图负责按从右往左的顺序把它们乘起来——顺序的选择就是反向模式与前向模式的分界，而反向要用前向留下的中间量这一条，正是显存账单的来源。

主线在图的最下面那一行：手推一次不是为了以后都手推。是为了框架给出一个数时，你能判断它算的是不是你要的那个量，以及它算错时该从哪一行开始看。

\subsection{怎么读这一章}

这一章适合拿纸逐行核对形状。先手推最小二乘与 logistic 回归，把 \(X^{\trans}\delta\) 这个骨架和那份对账清单练熟；再看 softmax 的 Jacobian 为何在交叉熵下整块消失，顺带把 \(\logsumexp\) 的数值实现看一遍；最后把同一过程放回计算图，反向传播就不再是一个需要单独记忆的算法，只是伴随量沿图反向传递。

每一篇读完可以问同一组问题：这个梯度的形状由什么定住，缩放系数出现在哪一个节点上，以及如果实现写错了，症状会以什么形式出现在损失曲线上。第三个问题最值得花时间，因为它决定了你能不能在没有参考实现的情况下自查。

\subsection{本章篇目}

以下篇目按概念依赖排列；若从具体问题进入，也可以先打开最接近当前困惑的一篇，再沿篇内链接回补。

\begin{itemize}
\tightlist
\item \hyperref[sec:10-Matrix-Calculus/04-Applications/01]{``从最小二乘到 logistic 回归梯度''}；
\item \hyperref[sec:10-Matrix-Calculus/04-Applications/02]{``\texorpdfstring{Softmax 交叉熵为何化简为 \(p-y\)}{Softmax 交叉熵为何化简为 p-y}''}；
\item \hyperref[sec:10-Matrix-Calculus/04-Applications/03]{``计算图上的反向传播与自动微分''}。
\end{itemize}

三篇走完，矩阵微积分这一部分的账就结清了：公式不必背，因为模型结构一变就能重新生成一遍。
